% Chapter file: 157_Zeitpfeil_En_ch.tex
% English Version
% Compatible with the English preamble (2026)

\section*{Abstract}
The fundamental equations of physics --- Newton, Maxwell, Schrödinger, Dirac, Einstein --- are all time-reversal invariant. Replacing $t$ with $-t$ leaves the equations valid. Why, then, do we experience a directed time? Standard physics answers this exclusively through the second law of thermodynamics: entropy increases. But this is a statistical law, not a fundamental explanation. T0 theory offers a geometric resolution: the time-mass duality $T \cdot m = 1$, the torus geometry of the vacuum field, and the fractal spacetime structure break time-reversal symmetry at the fundamental level.

\section{Why the Question Arises at All}

\subsection{Time-Reversal Invariance of the Fundamental Equations}

The deepest irritation in theoretical physics can be stated in one sentence: \emph{Not a single fundamental equation knows a preferred direction of time.}

Newton's second law $F = m\ddot{x}$ contains the second time derivative --- whether the film runs forward or backward, the equation remains the same. The same holds for the Maxwell equations, the Schrödinger equation, and the Einstein field equations. Even the Dirac equation, the relativistic quantum mechanics of the electron, is invariant under $t \to -t$ (combined with appropriate transformations).

Concretely: if $x(t)$ is a solution of $F = m\ddot{x}$, then $x(-t)$ is also a solution. A film showing a planetary orbit looks just as physical in reverse as it does forward. Two billiard balls colliding and flying apart --- the reverse process is equally permitted.

\subsection{The CPT Theorem}

Quantum field theory sharpens the problem. The CPT theorem states: every local, Lorentz-invariant quantum field theory with a Hermitian Hamiltonian is invariant under the combined transformation of charge conjugation~(C), parity inversion~(P), and time reversal~(T). Even the weak interaction, which violates C and P individually, respects CPT.

The observed CP violation in the kaon system does imply a T violation (to preserve CPT), but this is tiny ($\sim 10^{-3}$) and affects only specific decay processes. It does not explain why \emph{all} macroscopic processes possess a time direction.

\subsection{The Thermodynamic Way Out --- and Its Problem}

The only equation in classical physics with a built-in time direction is the second law of thermodynamics:
\begin{equation}
	\Delta S \geq 0
\end{equation}

Entropy increases --- eggs shatter but do not reassemble themselves. Yet Boltzmann showed that this law is \emph{statistical} in nature. It follows from probability, not from a fundamental asymmetry of the laws of nature. This merely shifts the question: why was the entropy of the universe initially low? Why does a ``Past Hypothesis'' (an initial condition of low entropy) exist?

Penrose quantified the problem: the initial state of the universe has a probability of $\sim 10^{-10^{123}}$ --- a number that seems to defy any physical explanation (cf.\ Doc.~030).

\subsection{Summary of the Problem}

\begin{itemize}
	\item All fundamental equations are $t \to -t$ invariant
	\item The CPT theorem guarantees this symmetry in QFT
	\item The observed CP/T violation is too weak for the macroscopic arrow of time
	\item The second law is statistical, not fundamental
	\item The low initial entropy remains unexplained
\end{itemize}

The question is therefore not: ``Why does time pass?'' --- but: ``Why does it pass in only \emph{one} direction, even though the equations permit both?''


\section{The T0 Answer: Geometric Time Direction}

T0 theory shifts the arrow of time from statistics to geometry. The time direction is not an emergent property of large particle numbers but a fundamental consequence of the spacetime structure.

\subsection{Argument 1: The Duality $T \cdot m = 1$ Is Asymmetric}

The central equation of T0 theory (cf.\ Doc.~069, 078) reads:
\begin{equation}
	T(x,t) \cdot m(x,t) = 1
\end{equation}

This equation looks symmetric --- but it is not. Mass~$m$ is \emph{positive definite}: there is no negative mass. This expressly includes antimatter --- a positron possesses exactly the same positive mass as an electron ($m_{e^+} = m_{e^-}$), only the opposite charge. The charge conjugation~C in CPT reverses charges, not the sign of mass. (This was experimentally confirmed at CERN in 2023: antihydrogen falls downward in a gravitational field, not upward.) It follows necessarily that the time field~$T$ is also positive definite: $T > 0$ everywhere and always.

In standard physics, the time coordinate~$t$ can take positive and negative values. In T0, the time field $T = 1/m$ is bound to mass and inherits its positivity. Time reversal $T \to -T$ would require $m \to -m$ --- physically impossible.

\emph{The duality $T \cdot m = 1$ breaks time-reversal symmetry algebraically.}

\subsection{Argument 2: The Torus Geometry Has a Winding Direction}

In T0 theory, the vacuum field is organized on a torus (cf.\ Doc.~018, 145). The phase~$\theta$ of the field winds along the torus:
\begin{equation}
	\dot{\theta} = m = \frac{1}{T}
\end{equation}

The winding has a \emph{chirality} --- a handedness. On a torus one can wind left or right, but not both simultaneously. Once the winding direction is chosen, the time direction is fixed.

This is analogous to a topological argument: a knot in a string has a handedness that cannot be reversed by continuous deformation. The winding of the vacuum field on the torus defines a topologically protected time direction.

\subsection{Argument 3: Fractal Spacetime Breaks T-Symmetry}

The local spacetime dimension in T0 deviates minimally from~3 (cf.\ Doc.~133):
\begin{equation}
	D_f^{\,\text{space}} = 3 - \xi \approx 2.9999
\end{equation}

A perfectly three-dimensional spacetime possesses perfect mirror symmetry --- including time reversal. But a fractal spacetime with $D_f < 3$ lacks perfect self-similarity under reflection. The ``porosity'' of spacetime is direction-dependent.

Formally: in a fractal space, path integrals for forward-running and backward-running paths differ. The number of available paths is direction-dependent because the fractal structure converts the discrete symmetry $t \to -t$ into a continuous asymmetry.

\subsection{Argument 4: Emergence of Time from Symmetry Breaking}

In the T0 model, time does not exist fundamentally but emerges from the $\xi$-field asymmetry (cf.\ Doc.~078). The hierarchy reads:
\begin{enumerate}
	\item Timeless, symmetric universe
	\item $\xi$-asymmetry arises (symmetry breaking)
	\item Time-mass duality $T \cdot m = 1$ becomes active
	\item Directed time stabilizes
\end{enumerate}

The symmetry breaking chooses a direction --- just as a pencil balanced on its tip chooses a direction when it falls. The chosen direction is then topologically protected by the torus winding.

\section{Why the Thermodynamic Arrow of Time Is Secondary}

In standard physics, the thermodynamic arrow of time is the \emph{only} argument for time direction. In T0, it is a \emph{consequence}, not the cause.

The logic is reversed:
\begin{enumerate}
	\item The geometry (torus, fractal dimension, duality) defines the time direction
	\item In a geometrically directed time, more microstates exist in the forward direction
	\item From this follows $\Delta S \geq 0$ as a geometric consequence
	\item The ``Past Hypothesis'' becomes unnecessary --- the low initial entropy is not an initial condition but a geometric necessity
\end{enumerate}

The fractal spacetime with $D_f < 3$ acts like a ratchet mechanism: path integrals have more available paths in one direction than in the other. Entropy increases \emph{because} the geometry possesses a preferred direction --- not the other way around.

\section{Scale Dependence of Time Direction}

The time direction is not equally pronounced at all scales (cf.\ Doc.~078 for details):

Below the scale $L_0 = \xi \cdot \ell_P$, time is granulated and chaotic --- here no stable time direction exists. In the transition zone around~$L_0$, the duality $T \cdot m = 1$ becomes active and the time direction begins to stabilize. Above~$L_0$, time is continuous and the direction is stable.

This explains why quantum mechanical equations \emph{appear} time-reversal invariant: at the quantum scale, the geometric asymmetry is so small ($\sim \xi \approx 10^{-4}$) that it is invisible in the equations. Only over many orders of magnitude does the effect accumulate --- analogous to the fractal correction $K_{\text{frak}} = 1 - 100\xi$ (cf.\ Doc.~133).

\section{Comparison with Other Approaches}

\begin{table}[htbp]
	\centering
	\begin{tabular}{p{3.5cm}p{4.5cm}p{4.5cm}}
		\toprule
		\textbf{Approach} & \textbf{Explanation} & \textbf{Weakness} \\
		\midrule
		Boltzmann & Statistical entropy increase & Does not explain low initial entropy \\
		Penrose (Weyl) & Weyl curvature hypothesis & Additional postulate required \\
		Carroll/Chen & Multiverse with spontaneous inflation & Not falsifiable \\
		\textbf{T0 theory} & \textbf{Geometric: torus + fractal spacetime + duality} & \textbf{Dependent on T0 framework} \\
		\bottomrule
	\end{tabular}
	\caption{Comparison of different explanations of the arrow of time}
\end{table}

\section{Summary}

The question ``Why is time directed?'' is a fundamental puzzle in standard physics because all fundamental equations are time-reversal invariant. The second law of thermodynamics does provide a time direction but is statistical and displaces the problem onto the unexplained low initial entropy.

T0 theory resolves this problem geometrically. Four arguments interlock: the duality $T \cdot m = 1$ forces $T > 0$ algebraically. The torus geometry defines a topologically protected winding direction. The fractal spacetime with $D_f < 3$ breaks mirror symmetry. And the emergence of time from symmetry breaking irreversibly selects a direction.

The thermodynamic arrow of time thereby shifts from cause to consequence: entropy increases \emph{because} the geometry is directed --- not the other way around.

\begin{thebibliography}{99}
	\bibitem{t0_zeit}
	Pascher, J., \emph{T0 Theory: Time as an Emergent Field}, Document~078.
	
	\bibitem{t0_zeitkonstant}
	Pascher, J., \emph{T0 Theory: Time and Energy in the T0 Model}, Document~069.
	
	\bibitem{t0_fraktal}
	Pascher, J., \emph{T0 Theory: Fractal Correction --- Derivation}, Document~133.
	
	\bibitem{t0_torus}
	Pascher, J., \emph{T0 Theory: Anomalous Magnetic Moments --- Torus Geometry}, Document~018.
	
	\bibitem{t0_ffgft}
	Pascher, J., \emph{FFGFT: Field Geometry on the Torus}, Document~145.
	
	\bibitem{t0_penrose}
	Pascher, J., \emph{T0 Theory and Penrose: Spacetime Structure}, Document~030.
	
	\bibitem{penrose_road}
	Penrose, R., \emph{The Road to Reality}, Jonathan Cape, 2004.
	
	\bibitem{boltzmann}
	Boltzmann, L., \emph{Vorlesungen \"uber Gastheorie}, Barth, 1896.
	
	\bibitem{carroll_eternity}
	Carroll, S., \emph{From Eternity to Here}, Dutton, 2010.
\end{thebibliography}
